\section{Многоуровневая модель атаки}

На основе \cite{latypov2020multilevel} выделяются три уровня абстракции, отражающие иерархию целей и средств злоумышленника.

\subsection{Стратегический уровень}

На этом уровне формулируются \textbf{бизнес-цели} атаки:
- Кража персональных данных (в соответствии с 152-ФЗ);
- Нарушение целостности финансовых транзакций;
- Вывод критической инфраструктуры из строя.

Каждая цель представляется как вершина с атрибутами \texttt{goal\_type}, \texttt{impact\_level}, \texttt{regulatory\_requirement}.

\subsection{Тактический уровень}

Тактики определяют \textbf{категории поведения} атакующего. Используется таксономия MITRE ATT\&CK, включающая 14 тактик, от Initial Access до Impact. Каждая тактика — вершина с атрибутами \texttt{tactic\_id}, \texttt{description}, \texttt{phase}.

\subsection{Методический уровень}

На этом уровне описываются \textbf{конкретные техники}:
- T1566.001 — Spearphishing Attachment;
- T1059.001 — PowerShell;
- T1071.001 — Web Protocols.

Вершины снабжаются атрибутами \texttt{technique\_id}, \texttt{cve}, \texttt{tool}, \texttt{detection\_rule}.

\subsection{Межуровневые связи и их семантика}

Связи между уровнями не являются произвольными. Они отражают причинно-следственные и логические зависимости:
- Дуга \texttt{realizes(g, t)} означает, что достижение цели \(g\) требует применения тактики \(t\);
- Дуга \texttt{implements(t, m)} указывает, что тактика \(t\) может быть реализована через метод \(m\);
- Дуга \texttt{requires(m1, m2)} фиксирует зависимость: метод \(m1\) невозможен без предварительного выполнения \(m2\).

Эти связи формируют \textbf{иерархический атрибутивный метаграф}, который служит основой для построения эталонных паттернов атак.